\chapter{Methodology} \label{chap:methods} The recent scientific revolution in data science, acquisition, and pre-processing has become readily accessible to those outside the strict doctrines of computer science through unique online learning programs, literature, and university classrooms. Institutions, groups, and labs such as the EPA monitoring service, Tropospheric Ozone Assessment Report (TOAR), and Global Monitoring Laboratory at the National Oceanic and Atmospheric Administration (GML-NOAA) have all provided easy access to surface level data and monitoring networks (U. EPA 2013; Schultz et al. 2017). Many researchers have utilized these extensive data stores and machine learning/artificial intelligence ensembles to assign exposures and deduce health trends (\cite{abdullah_development_2019}; \cite{chen_association_2023}; \cite{xue_hourly_2022}). Large databases generated by various disciplines utilizing established GDAL libraries and ESRI suite programs for the analysis of air pollutants, land use classification, and economic strategies have offered limitless development options for this project (\cite{rouault_gdal_2025}; \cite{dangermond_building_2020}). 

As technologies improve, advanced analytical methods become increasingly accessible, emerging continually in current and future disciplines (\cite{miller_data-driven_2015}). Machine learning (ML) and artificial intelligence (AI) models yield optimal results when pre-existing conditions and assumptions about the dataset are satisfied, such as non-stochasticity and a normalized distribution of training features, as outlined in \textcite{keller_machine_2017}. Extensive preprocessing and separation of training features (e.g., \cite{ilic_assessment_2020};\cite{napi_development_2021};\cite{staehle_quantifying_2022}) were undertaken due to the covariance between meteorological features and constituents that affect surface O$_3$, as discussed in Chapter \ref{chap:litreview}. Each pixel was defined by a complex ensemble and assigned a value, with its corresponding uncertainty estimated from an RK trend assessment of the same features, hereafter referred to as $\beta$. Values for $\beta$ and additional data sources across available time periods are presented in Chapter \ref{tmas}. The final selection of variables for each model and time frame used in this project is documented at the end of Chapter \ref{tmas:findatasets}.

This chapter outlines the implementation of specific CTM constructs in $\beta$ and demonstrates how surface O$_3$ is characterized in PHOTUC by predicting the residual geospatial error at each location. The prediction is applied to pixels within a raster brick created in Python. Discrepancies between the predictive features of surface O$_3$ and their true values are minimized through $\beta$’s use of CTMs, emissions-based features, and the geospatial uncertainty of $\beta$ at monitoring locations. This thesis combines similarities in spatial-temporal cycles, like-minded sciences, and the three laws of geography to propose a novel geospatial method for incorporating geospatial uncertainty into ML/AI ensembles (\cite{cao_deep_2022}; \cite{liu_improve_2018}). Each of the proposed transformations in this thesis has been observed in various geographical studies utilizing similar numerical models with considerable success.
%%%%%%%%%%%%%%%%%%%%%%%%%%%%%%%%%%%%%%%%%%%%%%%%%%%%%%%%%%%%%%%%%%%%%%%%%%%%%%%%%%%%%%%%%%%%%%%%%%%%%%%%%%%%%%%%%%%%
\section{Pre-Processing} \label{methods:ppro} These can be observed in Figure \ref{fig:varsot}, displaying increased seasonal tendencies and variation by county within the AOI. These aspects are further analyzed and compared in Chapters \ref{tmas} and \ref{rlts}. Temporal ranges of satellite technologies were utilized to create different datasets for testing. For example, GRIDMet data has been available since January 1st, 1979, and is continuously updated until a couple of days before the present, as provided by ECMWF \cite{european_centre_for_medium-range_weather_forecasts_era5_2019}, facilitating the creation of the historical dataset mentioned in \ref{tmas:hd}. 
\begin{figure}[h!]
    \centering
    \includegraphics[width=0.85\linewidth]{ThesisFilePkg//figs/feats_ot_v3_2019_2024.png}
    \caption{Model Variables Over Time.}
    \label{fig:varsot}
\end{figure}
This thesis covers the differences in data available from 2019-2024 and displays daily predictions for 5 randomly chosen months across the time frame, resulting in about 150 images. Full applications of the code will have more than ten thousand images depicting different AOIs across a 20-to-40-year time span. A Python script performed the pulling requests, JSON processing, and conversions to a common tabular format for further cleaning (\cite{rossum_python_2010}). Monitoring data was filtered to include only the counties of Maricopa, Pinal, and Pima. This data served as point features for raster statistics and was saved in a comma-separated values (CSV) file, forming the foundation for the datasets presented at the end of Chapter \ref{tmas:findatasets}. 
%%%%%%%%%%%%%%%%%%%%%%%%%%%%%%%%%%%%%%%%%%%%%%%%%%%%%%%%%%%%%%%%%%%%%%%%%%%%%%%%%%%%%%%%%%%%%%%%%%%%%%%%%%%%%%%%%%%%
\section{Vector Locations and Spatial Corrections} \label{methods:vlsc} The AZ State shapefile (\cite{united_states_census_bureau_tigerline_2022}) was downloaded at a census tract resolution for the relevant years within the specified timeframe. Unique Identifiers (UID) were composed of state, county, and census tract codes, which were spatially extracted in Python. ArcGIS PRO (ArcPRO) was later utilized due to its extreme presence in geographic academia and technical settings. Python facilitated the rapid cleaning, duplication, and distribution of project materials, while ArcGIS Pro enabled high-quality enhancement of the created data. Both programs performed spatial-joins for the TIGER-Line shapefile and raster representations of PHOTUC. The shapefile was also used for the creation of a mask for use in raster extraction, display, and overall uniformity of features during training. This mask was initially created at a 30 m resolution from a DEM provided by the U.S.G.S (\cite{us_geological_survey_1_2024}) and was resampled to 250 m for the final images.

The US Census Bureau contains extensive statistical collections of housing and income data relevant to the study area. This data was utilized to create visually appealing bivariate maps of the Area of Interest (AOI) seen later in Chapter \ref{sec:colordesc} and Appendix \ref{appendixA}. 28 were utilized, with an additional 65 monitors available for use in Arizona. Overall, the United States has thousands of monitors for use in future work. Maricopa, Pima, and Pinal correspond to county codes 013, 019, 021 respectively (United States Census Bureau 2025). Chapter II revealed that O$_3$ mechanisms are heavily based on energy components, urban presence, and meteorological drivers (A. M. Abdullah et al. 2017; Ye, Wang, and Zhang 2022) which are common in these three counties and Arizona as a whole. Some features were used to create unique predictor variables for use in the ML/AI ensembles based on CTM and Physics concepts learned in the researcher’s early academic career. 
%%%%%%%%%%%%%%%%%%%%%%%%%%%%%%%%%%%%%%%%%%%%%%%%%%%%%%%%%%%%%%%%%%%%%%%%%%%%%%%%%%%%%%%%%%%%%%%%%%%%%%%%%%%%%%%%%%%%
\section{Dependent and Independent Variable Extraction} \label{methods:varextraction} Necessary corrections, transformations, and resampling methods were made to reduce the complexity of error and implement estimated trends into the final models. The assumed independent variable, daily maximum 8-Hour surface O$_3$ detections, was selected based on extensive scientific evidence regarding the effects of surface ozone concentrations on public health and transport trends on the surface (Travis and Jacob 2019; M. Wang et al. 2016; W. Zhang et al. 2022); assumed potential dependent variables such and surface temperatures and precipitation were similarly gathered from known physical and chemical factors. Due to collaboration efforts among geospatial engineers, filtering is based on values commonly found in conventional GIS databases. The Air Quality Systems (AQS) \parencite{epa_air_2025} Application Programming Interface (API) was accessed using the "dailyData” and “byState" endpoints via in Python. Monitor locations whose center lied within the PHOTUC region were gathered, filtered, and geospatially corrected based on their provided coordinate reference systems (CRS). A data was re-projected to a meter (m) based projection based on the European Petroleum Spatial Grid (EPSG) identifier 32612 using data provided by epsg.io (Pridal et al. 2014). Some monitors before 2010 were reported with the North American Datum 1983 (NAD83; EPSG:4269) and were transformed to the World Geodesic System 1984 (WGS84; EPSG:4326) for spatial congruency.
%%%%%%%%%%%%%%%%%%%%%%%%%%%%%%%%%%%%%%%%%%%%%%%%%%%%%%%%%%%%%%%%%%%%%%%%%%%%%%%%%%%%%%%%%%%%%%%%%%%%%%%%%%%%%%%%%%%%
\section{Transformations And Theory Based Features} \label{methods:featsfm} Domain-inspired features were calculated to determine if depictions based on physics and chemistry could enhance the model without fully employing a CTM model. Cloud base and top height differentials were combined with pressure gradients and overall cloud fraction to represent potential cloud energy radiating within the atmosphere. Weekly moving averages were applied to multiple variables to capture temporal smoothing and remove high-frequency noise. Each feature implemented in the final model reduced spurious correlation between co-variates by combining their covariance into a representation of O$_3$ based on transport and chemical mechanisms. The second law of classical mechanics by Issac Newton (F=mass*acceleration) can be applied to bonds between molecules (\cite{taylor_classical_2005}); it is equal to the amount of heat an object generates over some distance, known as work ($W=Fdx$; \cite{borgnakke_fundamentals_2014}). This section details the overall extraction process, feature transformations, and thermodynamics of chemical states used in the final algorithm. 
%%%%%%%%%%%%%%%%%%%%%%%%%%%%%%%%%%%%%%%%%%%%%%%%%%%%%%%%%%
\subsection{Toms/Omi Dobson Unit (Du) Conversion to Surface Representation} \label{methods:tomi} One Dobson Unit is the number of molecules of ozone that would be required to create a layer of pure ozone 0.01 millimeters thick at a temperature of 0 degrees Celsius and a pressure of 1 atmosphere. Ground-based O$_3$ measurements taken from the EPA monitors are represented in parts-per-million (ppm). As stated in the Data Sources section, the TOMS/OMI satellite total column estimates are in Dobson Units (DU). 1 Dobson Unit would contain about $2.69x10^{16}$ ozone molecules for every square centimeter of area at the base of the column. For an array O$_3$ values in DU where 1 mole of Ozone is 47.9982 grams:
\[
1\,DU(O_3) = 2.69 \times 10^{16} \frac{O_3\,(\text{mol})}{\text{cm}^2} 
\times \frac{47.9982\,O_3\,(\text{g})}{6.0221 \times 10^{23}\,O_3\,(\text{mol})} 
\times \frac{1\,\text{mg}}{10^6\,\text{g}} 
\times \left(\frac{100\,\text{cm}}{1\,\text{m}}\right)^2 
\]

\[
=\frac{1.2912\times{10}^{18}\ mg\left(O_3\right)}{{6.0221\times{10}^{29}\ cm}^2}\ast\frac{1\times{10}^4cm^2}{1m^2}
\]

\[
=\frac{1.2912\times{10}^{22}\ mg\left(O_3\right)}{6.0221\times{10}^{29}m^2}
\]

\[
1\,DU(O_3) = 2.1441 \times 10^{-7} \frac{\text{mg}(O_3)}{\text{m}^2}
\]
As parts-per-million (ppm) is a dimensionless quantity used to relate weight to volume, this rendition of DU requires some form of divisible meter metric to be an accurate representation of surface ozone at some location. It’s quite convenient as a geographer to have access to a representative surface that can be stated in terms of x, y, and z. If one liter (L) can be represented as 0.0001m3, the ppm representation of the Dobson Unit per m$^2$ can be stated as:
\[
ppm(O_3, x, y) 
= \frac{mg(O_3)}{L} 
= \frac{mg(O_3)}{0.001\,\text{m}^3} 
= \frac{2.1441 \times 10^{-7}\,mg(O_3)}{0.1\,\text{m} \cdot \text{m}^2}
\]

\[
ppm(O_3, x, y) = 2.1441 \times 10^{-4} \,\hat{X}(DU)
\]
Where the vector $\hat{\mathbf{X}}$ is the monitor value of TCO in Dobson Units. Using this transformation, the TCO measurement from TOMS/OMI was converted to a surface representation of ozone in $\frac{mg\left(O_3\right)}{m^2}$. This was used to scale the DU down to a comparable measurement with Sentinel 5-P Imagery.
%%%%%%%%%%%%%%%%%%%%%%%%%%%%%%%%%%%%%%%%%%%%%%%%%%%%%%%%%%
\subsection{Kinetic Energy Representations} \label{methods:ke} Briefly switching into the realm of chemists and their amazing ability to provide concrete evidence of physical processes without the mathematical relationship via a tool called observation, chemical reactions always happen at a given rate. Calculated temperatures, concentrations, and unique combinations of constituents can be observed to do many different things, such as ice melting to water, then evaporating to vapor. Chemically, this is all H2O, simply in different thermodynamical states of a solid, liquid, and gas. As simplistically as possible (assuming no container (volume (V(L)(fluid + gas)) = 0), this can be mathematically represented as:
\[
Q(E) = mc_s \Delta T_{\text{ice}} + mc_s \Delta T_{\text{liquid}} + mc_s \Delta T_{\text{gas}}
\]
As mentioned in Chapter II.2, surface ozone reacts with certain VOC’s and emissions, oxidizing organic and inorganic material and causing what we know as surface O$_3$. Chapter II.3.1 showed that the remaining parts of this pseudo-catalyst are typically preserved and reintegrated into the atmosphere if given enough energy overtime. This has many different chemical representations; however, this thesis mainly focuses on nitrous oxides and carbonous aerosols as these are two of the most chemically influential chemicals of surface O$_3$ production (Cheng et al. 2018; Li et al. 2024; Richter 2009). Chapter II.3.3 revealed that carbon and nitrous based VoC emissions impact surface-level air quality, tropospheric photochemistry, and climate feedback mechanisms in addition to sustaining surface ozone concentrations:

Chapter 2 revealed that carbon and nitrous based VoC emissions impact surface-level air quality, tropospheric photochemistry, and climate feedback mechanisms in addition to sustaining surface O$_3$ reactions. These chemical relations are depictions of the sub-photosynthesis process, photolysis which has been found to drive O$_3$ (C. Chen et al. 2021; He and Carmichael 1999; Shah et al. 2024). In more thermodynamical terms, it is an energy-based oxidation event. If a constituent is required to produce O$_3$ and UV rays are necessary for ozone production, then it must exist in a constant reactionary state. The first law of thermodynamics states that matter can neither be created nor destroyed. This can be seen with the remaining O- that tend to titrate with VoCs. Atmospheric cycles attempt to rebalance the amount of newly seen negative oxygen molecules, in which energy is lost. This energy is measured by way of remote sensing and can theoretically be represented through a series of proper transformations. 
%%%%%%%%%%%%%%%%%%%%%%%%%%%%%%%%%%%%%%%%%%%%%%%%%%%%%%%%%%
\subsection{A Basic Implementation of Thermodynamics} \label{methods:thermo} O$_3$ is comprised of three oxygen (O$^{-}$) molecules. Known as a triatomic molecule, if one mol of oxygen is roughly 15.999 grams (g), then 1 mol of O$_3$ is about 47.9982 g. Given that air pollution operates within a system of finite degrees of freedom, the estimated kinetic energy must reflect the chemical and physical properties of its constituents. The essential conclusion is that the theoretical kinetic energies of any reaction at a specific state, in a defined space and time, can be expressed as a combination of chemical thermodynamics and derivations of the Ideal Gas Law. For subsequent use with a scaled field related to the predictor variable, the kinetic energy for a geo-chemical reaction could be represented as a function of spatial pressure, temperature, and volume dynamics of the system. The proposed average translational kinetic energy for a geo-atom z(x,y) is started as:
\[
{KE(x,y)}_t=\frac{1}{2}mv^2=\frac{3}{2}k_bT(x,y)
\]
Where at some time t, the estimated KE of a geo-atom can be represented as the and temperature $(T(x,y))$ of the constituent and kb, Boltzmann’s Constant. The short of this theory is that on average, atoms exist with a thermal energy of $\frac{3}{2}k_bT$, allowing the ideal gas law to be re-written as:
\[
KE\left(p,v,t\right)_{O_3}=\frac{3}{2}k_bT=DS\left(x,y\right)N\left(x,y\right)T(x,y)R
\]
Where ${P(x,y)}_t$ represents pressure, ${DS(x,y)}_t$ is used in place of the constant $\frac{3}{2}k_b$ as a measure of energy, R is the gas constant based on Boltzmann’s constant, and N(x,y) denotes the number of O$_3$ molecules estimated from remotely sensed imagery.
%%%%%%%%%%%%%%%%%%%%%%%%%%%%%%%%%%%%%%%%%%%%%%%%%%%%%%%%%%%%%%%%%%%%%%%%%%%%%%%%%%%%%%%%%%%%%%%%%%%%%%%%%%%%%%%%%%%%
\section{Interpolation and Imputation Strategies} \label{methods:inrtp} Linear, modified 2-Dimensional Akima, 1-Dimensional Akima, and polynomial functions of degrees 0, 1, 2, and 3 were evaluated using the same root mean square error (RMSE), mean absolute error (MAE), and mean square error (MSE) strategies as model evaluation methods. The creation of interpretation data for the effectiveness of interpolation methods is as follows:
\begin{itemize}
    \item Select all available monitor values
    \item Randomly changing 10 percent of these values to Nan in python
    \item Establish trend with missing monitors
    \item Calculate error statistics
    \item Place known values back into the dataset
    \item Predict with the established trend for unknown values
\end{itemize}
The $R^2$ correlation coefficient was added to interpolation-based methods, but not to regression-based imputation methods due it not being as valuable for low error in small predictions. A time-series based KNN imputation strategy was also developed for complex features missing in HD for further work on this project. Using the mean value between the best correlation methods, missing values were estimated and evaluated with methods in section \ref{methods:erroreval}. A Modified 2D Akima represents a moving piece-wise polynomial approximation over the given periods (\cite{akima_new_1970}). The trend estimated from training on replaced missing values was then used on the full data set, inclusive of the missing values taken out for KNN training. The trend derived from training on the dataset with replaced missing values was subsequently applied to the complete dataset, including the observations temporarily removed for KNN training. The error characteristics observed during interpolation were then used to adjust the corresponding missing rows in the final dataset, with the training-derived values replacing the initially interpolated or imputed estimates.
%%%%%%%%%%%%%%%%%%%%%%%%%%%%%%%%%%%%%%%%%%%%%%%%%%%%%%%%%%%%%%%%%%%%%%%%%%%%%%%%%%%%%%%%%%%%%%%%%%%%%%%%%%%%%%%%%%%%
\section{Missing Daily Raster Data} \label{methods:miar} Some datasets only provide yearly averages or estimations prior to 2000, due to sparse observation coverage and limited availability of consistent instruments and satellite records. Others face limitations stemming from the heterogeneity and metadata gaps typical of the large datasets used for routine monitoring. For example, nighttime light imagery is available with daily resolution only from 2013 onward; prior to this it must be accessed annually through another source as mentioned in Chapter \ref{tmas:ntl}. Normalized difference vegetation index (NDVI) imagery is subject to 16-day intervals to effectively “see” the change in wavelengths each source of vegetation is producing. To better account for daily changes, a function was created to calculate the difference between available temporal averages and daily differences to produce estimated daily changes. The function is simply a linear representation between the available data and number of days missing:
\[
dR\left(v,x,y\right)_\mu=\frac{\left({R(v,x,y)}_2-\ {R(v,x,y)}_1\right)}{N_{days}}
\]
Where $R\left(v,x,y\right)$ is a feature raster with a value, longitude, and latitude. $N_{days}$ is the difference in days from $t_1$ to $t_{max}$:
\[
N_{days}=t_{max}-t_1
\]
Once the average daily change was calculated, a new, estimated disaster comprised of the summation of the first date and the average daily change can be:
\[
\sum_{n=0}^{Days}{{R\left(v,x,y\right)}_{pred}={R\left(v,x,y\right)}_{1+n}+dR\left(v,x,y\right)_\mu}
\]
While this dataset may exhibit some repetitive values across monitor locations, this method facilitated the incorporation of a date-like index into the missing values.

For missing raster data, the time of each available rasters were extracted. The difference between the available raster data was divided by the number of missing days between the extracted times. The resulting mean was added to the first raster, producing a new raster for each day that was missing between the intervals. This had interesting effects on the outcome, working for features used in this these, but not so much for nighttime light imagery. While the imputation of new rasters from a preexisting set of models would have been ideal, this level of work would best be suited for a dissertation. This was only meant to introduce daily estimates into the algorithm to incorporate feature trends into each model.  
%%%%%%%%%%%%%%%%%%%%%%%%%%%%%%%%%%%%%%%%%%%%%%%%%%%%%%%%%%%%%%%%%%%%%%%%%%%%%%%%%%%%%%%%%%%%%%%%%%%%%%%%%%%%%%%%%%%%
\section{Statistical Model and Residual Kriging Methodology} \label{methods:smark} The Statistical Modeling and Residual Kriging (SMaRK) method is proposed as a methodology to leverage the flexibility of kriging alongside complex solutions from modern ensembles, utilizing satellite and monitoring systems for geo-pixels derived from satellite data. Chapter \ref{sec:sm} established that nonparametric, decision tree-based predictive models are increasingly adopted in spatial statistical analysis due to their adaptability in managing heterogeneous covariances (\cite{liu_improve_2018}; \cite{mu_estimation_2023}; \cite{pan_comparison_2023}). Many of the selected features during each training process showed covariance (i.e., specific humidity, temperature-created features, among others), ML/AI methods to minimize the influence of similar trends were required to define the overall property for each pixel. The SMaRK methodology is composed of five main steps:
\begin{enumerate}
    \item Statistically model a trend at known points given $\beta_n$, $\bm{f}\left(\bm{x}\left(\bm{s}\right);\hat{\beta}\right)$.
    \item Measure error at known locations given $f(\hat{\beta})$: $\varepsilon\left(s_0\right)$.
    \item Apply RK to estimate an error trend at given $\bm{x}\left(\bm{s}\right)$: $\bm{f}\left(\bm{x}\left(\bm{s}\right)\right) - \bm{x}\left(\bm{s};\hat{\beta}\right)$.
    \item Reincorporate estimated error into the established trend.
    \item Calculate the result as tupled pixels: $\langle\bm{x},Z_{sm}\left(\bm{x}\left(\bm{s}\right);\hat{\beta}\right),z\left(\varepsilon(\bm{s});\hat{\beta}\right)\rangle$.
\end{enumerate}
The resulting surface depicting O$_3$ properties was produced at a resolution of 250 m for each day within the selected time range. The underlying spatial concepts in SMaRK follow the establishment of initial geoatoms at monitoring locations based on predicted trends from a chosen ensemble. The full tuple of the geoatom from (\cite{goodchild_towards_2007}) in terms of statistical models and residual kriging to represent surface O$_3$ in PHOTUC on a given day can be depicted as:
\[
 \langle x,Z,z(x)t \rangle = \langle s_{0}(long;lat;t),Z(f(sm_{z};\beta(s_{0})),z{(s_0)}_{RK} \rangle
\]
The definition of property Z at a specific latitude, longitude, and time is given by $f_{sm}({\hat{z}}_{RK};\beta)$, which is derived from the estimated complex trend and the associated error for each monitoring location. A comprehensive theory based on a complex trend is refined by the residual laws linked to in situ measurements, resulting in $\hat{z}{(\mathbf{s})}_{RK}$, the estimated value of the system derived from $f(\hat{\beta})$. 
%%%%%%%%%%%%%%%%%%%%%%%%%%%%%%%%%%%%%%%%%%%%%%%%%%%%%%%%%%
\subsection{Statistical Parameters and Model Tuning} \label{methods:tuning} This thesis primarily discusses the sklearn package in Python, released by \cite{buitinck_api_2013}. Sklearn provides the most straightforward forms of modern complex analysis at the CPU processing level. The innovations within the sklearn development schemes have facilitated the integration of additional packages, serving as a foundation for implementing ML/AI modeling techniques into RK enhancements. XGBOOST is used as a separate entity and was able to be utilized with a GPU in the middle stages of this project. Although computation times are significantly increased for this thesis, further incorporation of these methods could lead to substantially reduced computation times and errors.

In this study, the models were trained using RandomSearchCV (RSCV) in Python. The search algorithm used a parameter list derived from previously known distributions of the data and information obtained in Chapter \ref{thermo}. The initial data set was reduced from approximately 2,760,000 to 656,000 rows in the 2018-2024 period. Keeping reproducibility in mind, each model was trained with a unique parameter combination, ensuring comparability to the data size and preventing overfitting or underfitting. After training, residual values were kriged using geospatial monitor positions and associated residual values, following methodologies from similar works (e.g, \cite{anand_gis-based_2021}; \cite{liu_improve_2018}; \cite{oliver_tutorial_2014}). A summary of the five numerical models and their potential parameters randomly selected in RSCV is as follows:
\begin{enumerate}
    \item \textbf{AdaBoost:} Learning rate \((10^{-5}\!-\!1)\) and the number of estimators \((50\!-\!500)\) were jointly tuned on \(N=\max(3,\lfloor p/2\rfloor)\) evenly spaced grids, with loss functions \{\emph{linear}, \emph{square}, \emph{exponential}\}. The best configuration minimized the cross‑validated validation error.
    \item \textbf{Gradient Boosting:} We searched over loss \{\emph{absolute error}, \emph{squared error}, \emph{Huber}\}, depth (None or \(3\!-\!10\)), estimators \([100,\lceil\max(500,10p)\rceil]\), learning rate \(10^{-3}\!-\!10^{-1}\) (log‑spaced), \( \alpha\in[0,0.1]\), and tolerance \(10^{-4}\!-\!10^{-2}\). Selection was by cross‑validated performance.
    \item \textbf{XGBoost:} We tuned estimators, learning rate \([0.001,\min(0.1,1/p)]\), grow policy \{\emph{depthwise}, \emph{lossguide}\}, and \(L_2\) regularization \(\lambda\in[0,1]\). The importance type was set for reporting. The chosen model balanced validation error and complexity.
    \item \textbf{Random Forest:} A grid over estimators, depth (None or \(3\!-\!10\)), split criterion, feature fraction \{\emph{log2}, \emph{sqrt}, \(0.75\), None\}, and \( \alpha\in\{0,0.001,0.05\}\) (post‑pruning) was evaluated with cross‑validation; the best forest maximized validation performance with stable generalization.
    \item \textbf{Multi‑Layer Perceptron:} Architectures from small (\((50)\), \((100)\)) to \(p\)-scaled stacks (\((p,2p,p)\), \((10p,5p,2p,p,1)\), etc.) were explored with activations \{\emph{logistic}, \emph{tanh}, \emph{ReLU}\}, \(L_2\) penalty \(10^{-6}\!-\!10^{-2}\), iteration caps, batch sizes \{\emph{auto}, \(p\), \(10p\)\}, and (when using Adam) \(\beta_1,\beta_2,\epsilon\) and learning‑rate schedule. The final network was the smallest that achieved the best cross‑validated score.
\end{enumerate}
Each of the models were set with a random state of 42. This number, as referenced in Douglas Adams' book "A Hitchhiker's Guide to the Galaxy" (1980) as the answer to the ultimate question of life, the universe, and everything, is used to "seed" results for this thesis and future modifications to the code. User-modified values for the features yield different results even with the same seeded values for each parameter; therefore, this option should be adjusted according to the area of interest (AOI) used. Additionally, the number of CPU cores utilized was set to the maximum number minus one to permit system operation during model training. After tuning, the in situ values from the monitoring locations were used to calculate a residual, which was predicted using a geospatial kriging of the same parameters. The more complex the ensemble, the more simplistic the error trends become, particularly in densely packed areas. 
%%%%%%%%%%%%%%%%%%%%%%%%%%%%%%%%%%%%%%%%%%%%%%%%%%%%%%%%%%%%
\subsection{Residual Krige Tuning} \label{methods:rkt} In geo-statistics (\cite{oliver_tutorial_2014}), regression kriging (RK) (\cite{chiles_fifty_2018}) is a hybrid approach for spatial prediction that combines statistical regression of independent variables based on dependent variables for trend estimation with a kriging term of the regression residuals to represent spatial dependence. If the field is defined by a representation of geoatoms, then this can be adapted to reflect the overall geo-field associated with its properties. For any location $s$ with a vector of observations $x(s)$, the RK estimator $\hat{z}{(\mathbf{s})}_{RK}$ is written as:
\[
\hat{z}{(\mathbf{s})}_{RK}=f\left(\mathbf{x}\left(\mathbf{s}\right);\hat{\beta}\right)+\varepsilon(\mathbf{s})
\]
Where $f\left(\mathbf{x}\left(\mathbf{s}\right);\hat{\beta}\right)$ is the regression term for the complex trend defined by the estimated parameters $\hat{\beta}$. The residual $\varepsilon\left(\mathbf{s}\right)$ is spatially auto-correlated and follows a selected variogram model. The RK method accommodates various regression approaches for base trend estimation, making it appropriate for utilizing complex trends estimated by ML/AI techniques. The parameters for potential variograms are as follows:
\begin{enumerate}
    \item{Variogram family:} linear, spherical, gaussian.
    \item{Pseudo-inverse:} False,True.
    \item{Pseudo-inverse type:} None, Hermitian pseudo‑inverse (pinvh).
    \item{Drift terms:} function vs. specified drift.
    \item{Number of lags:} 4,6,8.
    \item{Functional drift:} Fourier Series, Point Elevation, Sinusoidal.
\end{enumerate}
Each statistical model $f_{sm}(\hat{z}\left(\mathbf{s}\right);\beta)$ was fitted using the concentration of ground O$_3$ as the dependent variable. The geographic datasets and feature transformations served as predictors at each point. Given $V_1$ as the first feature used in the training, along with x remaining features:
\[
\hat{z}{(\mathbf{s}_\mathbf{0})}_{RK}=f_{sm}\left({{z}{(\mathbf{s}}_\mathbf{0})}_{v_1},{\mathbf{z}{(\mathbf{s}}_\mathbf{0})}_{v_2}\ldots,{\mathbf{z}{(\mathbf{s}}_\mathbf{0})}_{v_x};\hat{\beta}\right)+\varepsilon(\mathbf{s}_\mathbf{0})
\]
where every O$_3$ measurement at each monitoring station $z(\mathbf{s}_\mathbf{0})$ is defined by its corresponding residual $\varepsilon(\mathbf{s}_\mathbf{0})$, which is later compared to the associated estimation of the fitted set. This relationship facilitates interpolation by leveraging the inherent connections between the points, properties, and values of each monitor. The residual term, $\varepsilon\left(\mathbf{s}_\mathbf{0}\right)$, can be calculated through kriging estimated residuals at N monitoring stations:
\[
\varepsilon\left(\mathbf{s}_\mathbf{0}\right)=\left[\varepsilon\left(s_1\right),\ldots,\varepsilon\left(s_N\right)\right]^T=\lambda_0^T\varepsilon
\]
where the kriging weights $\lambda_0$ can be represented as a matrix composed of: 
\[
\lambda_0={\mathbf{C}(\beta)}^{-1}\mathbf{\hat{c}}_{\mathbf{0}}
\]
allowing for the covariance, $\mathbf{C}$, to be specified by $\beta_n$ among the O$_3$ monitor stations to account for spatial dependence. The covariance vector, $\mathbf{\hat{c}}_{\mathbf{0}}$ represents the spatial dependence between a monitor in the AOI, $\mathbf{s}_\mathbf{0}$, and its associated complex trend. Associations with $\mathbf{s}_\mathbf{0}$ can be further accounted for with an added drift term to covariance $\mathbf{C}$ by Universal Kriging (URK). 
%%%%%%%%%%%%%%%%%%%%%%%%%%%%%%%%%%%%%%%%%%%%%%%%%%%%%%%%%%%%
\subsection{Universal Kriging Drift Terms} \label{methods:urkt} Geostatistical literature employs various terms for techniques that are fundamentally similar. To avoid confusion, universal kriging (URK), kriging with external drift, and regression-kriging (RK) are the same technique (\cite{webster2007}). Originally termed Le krigeage universel by \textcite{matheron1969krigeage}, this technique was intended as a generalized approach in which trends are modeled as a function of coordinates. Consequently, many authors reserve the term universal kriging (UK) for situations where only the coordinates are used as predictors. When the deterministic component of variation (drift) is defined as a function of external auxiliary variables, RK becomes known as residual kriging with external drift (RKED) (\cite{hengl2004}). 

In the context of RKED, or URK, predictions are made similarly to standard kriging. However, the covariance matrix of residuals $\hat{z}{(\mathbf{s}_\mathbf{0})}_{RK}$ can be formulated using auxiliary predictors: 
\[\hat{z}{(\mathbf{s}_\mathbf{0})}_{RK}=\hat{z}^{\mathrm{URK}}(\mathbf{s}_0) = \sum_{i=1}^{n} w_i^{\mathrm{URK}}(\mathbf{s}_0) \, z(\mathbf{s}_i)\]
The drift and residuals can be estimated separately and summed, similar to methodology established in the 1990s with the advent of satellite technology (e.g., \cite{ahmed1987}; \cite{odeh1995}; \cite{journel1989}) wherein $k$ monitors can be further related to weighted terms known as spatial drift: 
\[
\hat{z}{(\mathbf{s}_\mathbf{0})}_{RK}=\sum_{i=1}^{n} w_i^{\mathrm{URK}}(\mathbf{s}_0) \cdot q_k(\mathbf{s}_i) =q_k(\mathbf{s}_0)
\] 
Replacing the weighted term, $w_i^{\mathrm{URK}}$, for each monitor with a transpose of the resulting estimation of $q_k(\mathbf{s}_0)$ denoted as $\boldsymbol{\delta}_0$:
\[
\hat{z}{(\mathbf{s}_\mathbf{0})}_{RK}=\hat{z}^{\mathrm{URK}}(\mathbf{s}_0) = \boldsymbol{\delta}_0^{T} \cdot \mathbf{z}
\] where $\mathbf{z}$ is the target residual, and $q_k$ denotes the predictor variables, i.e., values at a new location $\mathbf{s}_0$. The associated vector of URK weights can be derived from the number of predictors, and $\mathbf{z}$ as a vector of $n$ observations at primary locations. The weights, $w_i^{\mathrm{URK}}$, are computed using the extended matrices:
\[
\lambda_0^{\mathrm{URK}} =
\left\{ w_1^{\mathrm{URK}}(\mathbf{s}_0), \ldots, w_n^{\mathrm{URK}}(\mathbf{s}_0), \gamma_1(\mathbf{s}_0), \ldots, \gamma_p(\mathbf{s}_0) \right\}^{T}
\mathbf{C}^{\mathrm{URK-1}} \, \mathbf{c}_0^{\mathrm{URK}}
\]
Where $\lambda_0^{\mathrm{URK}}$ is the vector of solved weights for the Lagrange multipliers $\gamma_p$, $\mathbf{C}^{\mathrm{URK}}$ represents the extended covariance matrix of residuals, and $\mathbf{c}_0^{\mathrm{URK}}$ is the extended vector of covariances at each adjusted location. The full extended covariance matrix of residuals from \textcite{webster2007} can then be represented as:
\[
\mathbf{C}^{\mathrm{URK}} =
\begin{bmatrix}
C(\mathbf{s}_1, \mathbf{s}_1) & \cdots & C(\mathbf{s}_1, \mathbf{s}_n) & \gamma_1(\mathbf{s}_1) & \cdots & \gamma_p(\mathbf{s}_1) \\
\vdots & \ddots & \vdots & \vdots & \ddots & \vdots \\
C(\mathbf{s}_n, \mathbf{s}_1) & \cdots & C(\mathbf{s}_n, \mathbf{s}_n) & \gamma_1(\mathbf{s}_n) & \cdots & \gamma_p(\mathbf{s}_n) \\
\gamma_1(\mathbf{s}_1) & \cdots & \gamma_1(\mathbf{s}_n) & 0 & \cdots & 0 \\
\vdots & \ddots & \vdots & \vdots & \ddots & \vdots \\
\gamma_p(\mathbf{s}_1) & \cdots & \gamma_p(\mathbf{s}_n) & 0 & \cdots & 0
\end{bmatrix}
\] Where equations from generalized least squares combined with the unbiased kriging constraints are implemented through block-augmented matrices to reduce $\mathbf{c}_0^{\mathrm{URK}}$ to one:
\[
\mathbf{c}_0^{\mathrm{URK}} =
\left\{ C(\mathbf{s}_0, \mathbf{s}_1), \ldots, C(\mathbf{s}_0, \mathbf{s}_n), \gamma_0(\mathbf{s}_0), \gamma_1(\mathbf{s}_0), \ldots, \gamma_p(\mathbf{s}_0) \right\}^{T} ; \gamma_0(\mathbf{s}_0) = 1
\]
The drift functions for $\hat{z}^{\mathrm{URK}}(\mathbf{s}_0)$ included a sinusoidal estimation based on latitude and longitude, point-specific drift with elevation values, and a Fourier Series. Spatial drift was integrated with $C$ as a sinusoidal function and the elevation at the monitor location. Replacing $C$ with a Fourier series allows for the derivation of potential variations representing the magnitude of corresponding error at each monitor in URK. Given a specific wave function $f(x)$, a period of $2\pi$ can be expanded as 
\[f(x) = a_0 + \sum_{n=1}^\infty (a_n\cos nx +b_n\sin nx)\] 
where the coefficients of $f(x)$ are determined by Euler's formulas for each period: 
\[a_0 = \frac{1}{2\pi}\int_{-\pi}^\pi f(x)dx, a_n =  \frac{1}{\pi}\int_{-\pi}^\pi f(x)\cos (nx) dx,b_n =  \frac{1}{\pi}\int_{-\pi}^\pi f(x)\sin (nx) dx\]
for $n\geq 1$. Here, $f(x)$ represents the estimation derived from the complex trend $\varepsilon(\mathbf{s}_\mathbf{0}$. Mathematically, this establishes a direct relationship between complex trend estimations and their corresponding in situ values. This approach facilitates a more accurate representation of surface O$_3$ based on spatially dependent covariates. 
%%%%%%%%%%%%%%%%%%%%%%%%%%%%%%%%%%%%%%%%%%%%%%%%%%%%%%%%%%
\section{Model Evaluation Techniques} \label{methods:erroreval} The models were selected based on several factors. The mean absolute percent error (MAPE) and the root mean square error (RMSE) were primarily used for model comparison, as studies typically prefer metrics that are relative to the dataset (\cite{hu_optimized_2023}; \cite{mu_estimation_2023}; \cite{pan_comparison_2023}; \cite{zhang_deep_2023}). Additionally, the mean absolute error (MAE) and the correlation coefficient (R$^2$) were calculated. Although these metrics are not obligatory, this thesis also provides the minimums, maximums, medians, and overall distributions of the predicted data to offer deeper insights into the improvements achieved by incorporating geospatial uncertainty. The optimal model, whether it incorporates RK improvements or not, was applied to the full dataset for final raster predictions. 

RMSE and MAPE,, along with statistical distributions of the outcomes, serve as the primary determinants in the final model for Chapter \ref{chap:discussion}. Calculating R$^2$ is not typical in the context of nonlinear regression models, as values can sometimes occur outside the 0 to 1 range under specific conditions (\cite{mu_estimation_2023};\cite{zhang_deep_2023}). Due to RK improvements providing a simplified linear association with geospatial data, the interpretation of R$^2$ should still depict tendencies across each predicted value versus random events. Lower values typically indicate that the model can also apply to random associations within the data. However, it should not be relied upon if other metrics indicate low error. These error metrics were also employed for the evaluation of interpolation and the imputation of monitor values based on linear methods.
%%%%%%%%%%%%%%%%%%%%%%%%%%%%%%%%%%%%%%%%%%%%%%%%%%%%%%%%%%
\section{Conclusion} \label{methods:conclusion} All calculations were conducted in ppm and subsequently converted to ppb for improved clarity in the representation of differences. Models were evaluated using various combinations of the collected features to assess whether data from 1980 onwards can be leveraged to create historical models as effective as those based on contemporary data. The features were chosen based on their correlation with the daily surface O$_3$ values at the monitoring locations, existing literature, and insights gained through coursework. Initially, the cut-off points for the number of features included in each model were determined by the ranking of variables in each of the three correlation equations: Pearson, Spearman, and Kendall. 

Upon final review, the ranking of feature variables prior to their inclusion in the model guided the selection of the most pertinent variables for each data set presented throughout Chapter \ref{tmas:findatasets}. Only one data set used a cut-off point for the top 24 features in the Pearson correlations. A raster dataset was created at a spatial resolution of 500 m from imagery gathered in Google Earth Engine (GEE) using combinations of bilinear and cubic aggregation methods. The stacked result was resampled to 250 m with an RK approximation of the surface based on the best model. 

This was subsequently integrated into a 250 m resolution RK grid, developed by modeling uncertainty around the monitoring stations for the visualization of the final daily images. In total, five ML/AI ensembles were used to identify a trend that minimizes the complexity of error for establishing the RK trend. The overall predictive error was further reduced by incorporating the geospatial uncertainty of the selected features into the chosen ensemble. Even if the initial numerical model produced no error, the RK method addressed this and only enhanced regions further from the true value. 